\section{Setup}
\label{sec:setup}

\paragraph{Task, labels and kinds of error.}
A model receives a user request and a set of tool schemas and generates greedily. The call is parsed in whichever format the model writes and compared with the benchmark's ground truth after both are normalised to a name and an argument dictionary. Each call receives one failure mode: valid, wrong name, wrong argument values, missing arguments, missing parallel calls, an unneeded call, no call, an unparseable call, or a call cut off by the budget. Every number is computed on the scored population, the items where a call is required and the generation can be judged, and "every scored failure" means every failure mode in that population. We also score two kinds of error against the valid calls alone: wrong argument values, where the right tool gets a wrong value, and dropped parallel calls, where a request needs several calls and the model writes fewer.

\paragraph{Prompt route.}
Tool schemas reach the model through its own chat template where the template carries them. Gemma-3's template drops the tool argument, so every Gemma-3 run, and both forced-JSON runs, use a fallback system prompt that lists the schemas and asks the model to reply with only one JSON object of the form name and arguments and nothing else (\cref{app:protocol}). Rendering one prompt per model on CPU with the paper's render function confirms that Gemma-3 is the only one of the six whose template drops the schemas. A prompt that asks for one object invites a model to drop the other calls of a parallel request, so dropped calls on these runs are partly a property of the prompt, and \cref{sec:result} treats them that way.

\paragraph{Conditioning facts.}
The labelling and extraction code was audited and repaired before the measurements here (\cref{app:audit}). The largest repair concerns correct list-valued arguments, which had been scored wrong on the JSON formats: on Qwen3-1.7B it turned \repPosBeforeQwenThreeBfcl{} labelled failures into \repPosAfterQwenThreeBfcl{}. Relabelling the stored generations of \nRelabelRuns{} runs with the released labeller reproduces every stored label. Two extractor generations coexist (\cref{tab:provenance}): \nFirstExtractInternals{} of the \nInternalsRuns{} probe-side runs and \nFirstExtractConfidence{} of the \nConfidenceRuns{} confidence-side runs keep the first extraction, which records no commit, and the re-extracted runs do not all come from a clean tree. One run extracted under both reads \replGapStale{} and \replGapClean{}, a drift of \replGapDelta{}. Some failures are schema echoes, calls whose arguments copy the tool's JSON schema: a pattern for a copied \texttt{properties} field or object type finds them in \echoPosGemmaBfcl{}\% of Gemma-3 failures on BFCL and on no confidence-side run. Most are labelled missing arguments and stay out of the value-error comparison.

\paragraph{Models and benchmarks.}
The checkpoints are Llama-3.2-1B and 3B, Gemma-3-1B, Qwen3-1.7B, Qwen3.5-0.8B and MiniCPM5-2B, from four developers, released between September 2024 and September 2026 (\cref{app:models}). The chat templates of the Qwen3, Qwen3.5 and MiniCPM5 checkpoints carry a reasoning mode, which is disabled at generation throughout. The benchmarks are the Glaive function-calling corpus \citep{glaive2024}, the single-turn categories of BFCL-v4 \citep{patil2025bfcl} and its live categories, contributed by users. Glaive repeats requests: on Llama-3.2-1B its \gdNScoredLlamaOneBGlaive{} scored items hold \gdNUniquePromptsLlamaOneBGlaive{} unique requests. A run is one checkpoint on one benchmark. A run with fewer than thirty items of either class is underpowered and excluded from every aggregate, which leaves the confidence-side checkpoints with powered runs on BFCL and BFCL-live only. A kind of error is scored on a run only with at least twenty failures of that kind.

\paragraph{Detectors.}
Confidence is the mean token log-probability over the generated call, with its direction fixed on training items, chosen before any result was seen. It averages over every token of the call, while the probe reads the name, the argument values and the delimiter separately. The probe is the token-role probe of \citet{healy2026internal}, a regularised logistic regression on the residual state at the function name, the mean over argument-value tokens and the closing delimiter, at eight depths. The output-only judge is one logistic regression on what an operator sees from outside the model: lengths, benchmark category, counts read off the call, every stored confidence summary, and character n-grams of the call and the request. The reader is Qwen3.5-0.8B, which reads the request and the call as plain text, without schemas, and is probed with the same classifier. It is the same checkpoint as one confidence-side model. The floor is a classifier on prompt length, generation length and truncation.

\paragraph{Protocol and the two sides.}
Items are grouped by ground-truth tool and split five ways at the group level, so every score is on tools absent from training. Every trained detector is fitted on the training folds restricted to the scored population, with a validation carve-out of training tools. One pooled AUC is computed per split seed, and five seeds are averaged. Every difference between detectors is a paired bootstrap that resamples whole tools, with its 95\% percentile interval. The internal advantage of a run is the AUC of the probe minus that of confidence. Each run is assigned to a side by its family, fixed in the analysis code when the split was first described: Llama-3.2 and Gemma-3 on the probe side, Qwen3, Qwen3.5 and MiniCPM5 on the confidence side. The assignment came from the data it describes, and a run can disagree with its side.
